\PassOptionsToPackage{unicode=true}{hyperref} % options for packages loaded elsewhere
\PassOptionsToPackage{hyphens}{url}
%
\documentclass[14pt,]{extarticle}
\usepackage{lmodern}
\usepackage{amssymb,amsmath}
\usepackage{ifxetex,ifluatex}
\usepackage{fixltx2e} % provides \textsubscript
\ifnum 0\ifxetex 1\fi\ifluatex 1\fi=0 % if pdftex
  \usepackage[T1]{fontenc}
  \usepackage[utf8]{inputenc}
  \usepackage{textcomp} % provides euro and other symbols
\else % if luatex or xelatex
  \usepackage{unicode-math}
  \defaultfontfeatures{Ligatures=TeX,Scale=MatchLowercase}
    \setmainfont[]{Times New Roman}
  \ifxetex
    \usepackage{xeCJK}
    \setCJKmainfont[]{Microsoft JhengHei}
  \fi
  \ifluatex
    \usepackage[]{luatexja-fontspec}
    \setmainjfont[]{Microsoft JhengHei}
  \fi
\fi
% use upquote if available, for straight quotes in verbatim environments
\IfFileExists{upquote.sty}{\usepackage{upquote}}{}
% use microtype if available
\IfFileExists{microtype.sty}{%
\usepackage[]{microtype}
\UseMicrotypeSet[protrusion]{basicmath} % disable protrusion for tt fonts
}{}
\IfFileExists{parskip.sty}{%
\usepackage{parskip}
}{% else
\setlength{\parindent}{0pt}
\setlength{\parskip}{6pt plus 2pt minus 1pt}
}
\usepackage{hyperref}
\hypersetup{
            pdfborder={0 0 0},
            breaklinks=true}
\urlstyle{same}  % don't use monospace font for urls
\usepackage[margin=1in]{geometry}
\usepackage{longtable,booktabs}
% Fix footnotes in tables (requires footnote package)
\IfFileExists{footnote.sty}{\usepackage{footnote}\makesavenoteenv{longtable}}{}
\setlength{\emergencystretch}{3em}  % prevent overfull lines
\providecommand{\tightlist}{%
  \setlength{\itemsep}{0pt}\setlength{\parskip}{0pt}}
\setcounter{secnumdepth}{0}
% Redefines (sub)paragraphs to behave more like sections
\ifx\paragraph\undefined\else
\let\oldparagraph\paragraph
\renewcommand{\paragraph}[1]{\oldparagraph{#1}\mbox{}}
\fi
\ifx\subparagraph\undefined\else
\let\oldsubparagraph\subparagraph
\renewcommand{\subparagraph}[1]{\oldsubparagraph{#1}\mbox{}}
\fi

% set default figure placement to htbp
\makeatletter
\def\fps@figure{htbp}
\makeatother


\date{}

\begin{document}

\hypertarget{week-05-ux4f5cux696dux56deux994b}{%
\section{Week 05 作業回饋}\label{week-05-ux4f5cux696dux56deux994b}}

\begin{quote}
\textbf{評分標準：} 依據 \texttt{week05\_rubric\_v2.xlsx} 之 13
項評分指標，滿分 100 分。 各大項配分：A. 實驗架構（Coder）50 分 ／ B.
Builder 版本 20 分 ／ C. 程式品質 15 分 ／ D. 文件與繳交 15 分
\end{quote}

\begin{center}\rule{0.5\linewidth}{0.4pt}\end{center}

\hypertarget{ux5442ux6770ux9a5b-100100}{%
\subsection{113825002 呂杰驛 ---
100／100}\label{ux5442ux6770ux9a5b-100100}}

\hypertarget{ux7e3dux8a55}{%
\subsubsection{總評}\label{ux7e3dux8a55}}

架構完整度為本次作業最高，Coder 與 Builder 兩版本均完整實作繪圖、2AFC
判斷、資料儲存及結束畫面。多筆畫分離追蹤與座標幾何中心校正為其他同學未處理的進階細節，整體品質優秀。

\hypertarget{ux5404ux9805ux8a55ux8a9e}{%
\subsubsection{各項評語}\label{ux5404ux9805ux8a55ux8a9e}}

\begin{longtable}[]{@{}lll@{}}
\toprule
\begin{minipage}[b]{0.08\columnwidth}\raggedright
指標\strut
\end{minipage} & \begin{minipage}[b]{0.02\columnwidth}\raggedright
得分\strut
\end{minipage} & \begin{minipage}[b]{0.81\columnwidth}\raggedright
評語\strut
\end{minipage}\tabularnewline
\midrule
\endhead
\begin{minipage}[t]{0.08\columnwidth}\raggedright
A1 試驗迴圈\strut
\end{minipage} & \begin{minipage}[t]{0.02\columnwidth}\raggedright
15/15\strut
\end{minipage} & \begin{minipage}[t]{0.81\columnwidth}\raggedright
10 輪 \texttt{for\ trial\ in\ range(10)} 結構乾淨，每輪以 P
鍵啟動，節奏控制良好。\strut
\end{minipage}\tabularnewline
\begin{minipage}[t]{0.08\columnwidth}\raggedright
A2 刺激呈現與計時\strut
\end{minipage} & \begin{minipage}[t]{0.02\columnwidth}\raggedright
10/10\strut
\end{minipage} & \begin{minipage}[t]{0.81\columnwidth}\raggedright
使用 \texttt{core.CountdownTimer(10)}
控制繪圖階段並即時顯示剩餘秒數；2AFC 階段以 \texttt{CountdownTimer(5)}
限時，時間控制精準。\strut
\end{minipage}\tabularnewline
\begin{minipage}[t]{0.08\columnwidth}\raggedright
A3 反應收集與正確性\strut
\end{minipage} & \begin{minipage}[t]{0.02\columnwidth}\raggedright
15/15\strut
\end{minipage} & \begin{minipage}[t]{0.81\columnwidth}\raggedright
A／L 按鍵對應左右，配合 \texttt{is\_original\_left}
旗標判斷正確性，邏輯清晰無誤。\strut
\end{minipage}\tabularnewline
\begin{minipage}[t]{0.08\columnwidth}\raggedright
A4 區塊結構\strut
\end{minipage} & \begin{minipage}[t]{0.02\columnwidth}\raggedright
10/10\strut
\end{minipage} & \begin{minipage}[t]{0.81\columnwidth}\raggedright
每輪試驗前以「按 P
鍵開始」提示分隔，受試者可自行控制節奏；包含完整指導語畫面與結束畫面顯示總分。\strut
\end{minipage}\tabularnewline
\begin{minipage}[t]{0.08\columnwidth}\raggedright
B1 Builder 流程重建\strut
\end{minipage} & \begin{minipage}[t]{0.02\columnwidth}\raggedright
10/10\strut
\end{minipage} & \begin{minipage}[t]{0.81\columnwidth}\raggedright
完整重現 Coder 版本流程（Intro → Preparing → Drawing → 2AFC →
TaskOver），以 \texttt{trials\_loop} 控制 10 輪迴圈。\strut
\end{minipage}\tabularnewline
\begin{minipage}[t]{0.08\columnwidth}\raggedright
B2 Builder 資料紀錄\strut
\end{minipage} & \begin{minipage}[t]{0.02\columnwidth}\raggedright
5/5\strut
\end{minipage} & \begin{minipage}[t]{0.81\columnwidth}\raggedright
使用
\texttt{thisExp.addData(\textquotesingle{}Correct\textquotesingle{},\ is\_correct)}
將正確性記入 PsychoPy 原生資料檔，資料結構完整。\strut
\end{minipage}\tabularnewline
\begin{minipage}[t]{0.08\columnwidth}\raggedright
B3 Code Component / JS\strut
\end{minipage} & \begin{minipage}[t]{0.02\columnwidth}\raggedright
5/5\strut
\end{minipage} & \begin{minipage}[t]{0.81\columnwidth}\raggedright
Code Component 中提供 Auto→JS 對應版本，顯示已考慮 Pavlovia
線上部署的可行性。\strut
\end{minipage}\tabularnewline
\begin{minipage}[t]{0.08\columnwidth}\raggedright
C1 可讀性與組織\strut
\end{minipage} & \begin{minipage}[t]{0.02\columnwidth}\raggedright
5/5\strut
\end{minipage} & \begin{minipage}[t]{0.81\columnwidth}\raggedright
變數命名明確（\texttt{all\_strokes}、\texttt{is\_drawing}、\texttt{norm\_strokes}），邏輯分段清楚。小建議：\texttt{import\ pandas}
與 \texttt{import\ datetime} 目前放在檔案末端（第 161--162 行），Python
慣例建議將 import 集中在檔案開頭，此為風格建議，不扣分。\strut
\end{minipage}\tabularnewline
\begin{minipage}[t]{0.08\columnwidth}\raggedright
C2 Escape 處理\strut
\end{minipage} & \begin{minipage}[t]{0.02\columnwidth}\raggedright
5/5\strut
\end{minipage} & \begin{minipage}[t]{0.81\columnwidth}\raggedright
在準備畫面、繪圖階段、2AFC 階段均偵測 Escape 鍵，並以
\texttt{win.close();\ core.quit()} 正確清理資源後退出。\strut
\end{minipage}\tabularnewline
\begin{minipage}[t]{0.08\columnwidth}\raggedright
C3 邊界情況處理\strut
\end{minipage} & \begin{minipage}[t]{0.02\columnwidth}\raggedright
5/5\strut
\end{minipage} & \begin{minipage}[t]{0.81\columnwidth}\raggedright
當繪圖為空時以預設座標填充避免 ShapeStim
崩潰；座標歸零（幾何中心校正）確保翻轉後圖形居中對稱。\strut
\end{minipage}\tabularnewline
\begin{minipage}[t]{0.08\columnwidth}\raggedright
D1 繳交完整度\strut
\end{minipage} & \begin{minipage}[t]{0.02\columnwidth}\raggedright
5/5\strut
\end{minipage} & \begin{minipage}[t]{0.81\columnwidth}\raggedright
同時繳交 \texttt{.py} 與 \texttt{.psyexp}，附 \texttt{content.html}
詳細說明兩版本的輸出檔案格式與欄位意義。\strut
\end{minipage}\tabularnewline
\begin{minipage}[t]{0.08\columnwidth}\raggedright
D2 範例輸出\strut
\end{minipage} & \begin{minipage}[t]{0.02\columnwidth}\raggedright
5/5\strut
\end{minipage} & \begin{minipage}[t]{0.81\columnwidth}\raggedright
\texttt{content.html} 中清楚記載 CSV
欄位定義（Trial、Correct），可據以驗證資料品質。\strut
\end{minipage}\tabularnewline
\begin{minipage}[t]{0.08\columnwidth}\raggedright
D3 科學設計品質\strut
\end{minipage} & \begin{minipage}[t]{0.02\columnwidth}\raggedright
5/5\strut
\end{minipage} & \begin{minipage}[t]{0.81\columnwidth}\raggedright
左右隨機化（\texttt{random.choice}）、多筆畫分離追蹤（\texttt{all\_strokes}）、座標中心校正後再翻轉，以及
\texttt{try/except/finally}
結構確保資料儲存，實驗設計嚴謹且具巧思。\strut
\end{minipage}\tabularnewline
\bottomrule
\end{longtable}

\end{document}
